% LaTeX companion of hw07.typ. Both formats are editable.
\section{Homework 7 --- submitted work}

\chapter{Part A}

\section{4.3 Exercise 14}

For \(T(M) = \begin{pmatrix}
1 & 1 \\
2 & 2
\end{pmatrix}M\) relative to \(\mathcal{B} = \left( {\begin{pmatrix}
1 & 0 \\
{- 1} & 0
\end{pmatrix},\begin{pmatrix}
0 & 1 \\
0 & {- 1}
\end{pmatrix},\begin{pmatrix}
1 & 0 \\
2 & 0
\end{pmatrix},\begin{pmatrix}
0 & 1 \\
0 & 2
\end{pmatrix}} \right)\), the submission forms \(\lbrack T\rbrack_{\mathcal{B}}\) from \(\left\lbrack {T\left( b_{i} \right)} \right\rbrack_{\mathcal{B}}\). It finds

\(T\left( b_{1} \right) = T\left( b_{2} \right) = \begin{pmatrix}
0 & 0 \\
0 & 0
\end{pmatrix},\quad T\left( b_{3} \right) = \begin{pmatrix}
3 & 0 \\
6 & 0
\end{pmatrix},\quad T\left( b_{4} \right) = \begin{pmatrix}
0 & 3 \\
0 & 6
\end{pmatrix},\)

so

\(\lbrack T\rbrack_{\mathcal{B}} = \begin{pmatrix}
0 & 0 & 0 & 0 \\
0 & 0 & 0 & 0 \\
0 & 0 & 3 & 0 \\
0 & 0 & 0 & 3
\end{pmatrix}\mapsto\begin{pmatrix}
0 & 0 & 1 & 0 \\
0 & 0 & 0 & 1 \\
0 & 0 & 0 & 0 \\
0 & 0 & 0 & 0
\end{pmatrix}.\)

Hence

\(\left\lbrack {\text{ker}\ T} \right\rbrack_{M} = \text{span}\left( {\begin{pmatrix}
1 \\
0 \\
0 \\
0
\end{pmatrix},\begin{pmatrix}
0 \\
1 \\
0 \\
0
\end{pmatrix}} \right),\)

\(\left\lbrack {\text{im}\ T} \right\rbrack_{M} = \text{span}\left( {\begin{pmatrix}
0 \\
0 \\
1 \\
0
\end{pmatrix},\begin{pmatrix}
0 \\
0 \\
0 \\
1
\end{pmatrix}} \right).\)

A basis for the kernel is \(\left( {\begin{pmatrix}
1 & 0 \\
{- 1} & 0
\end{pmatrix},\begin{pmatrix}
0 & 1 \\
0 & {- 1}
\end{pmatrix}} \right)\); a basis for the image is \(\left( {\begin{pmatrix}
1 & 0 \\
2 & 0
\end{pmatrix},\begin{pmatrix}
0 & 1 \\
0 & 2
\end{pmatrix}} \right)\); and \(\text{rank}(T) = 2\).

\section{4.3 Exercise 28}

For \(T\left( {f(t)} \right) = f\left( {2t - 1} \right)\) and \(\mathcal{B} = \left( {1,t - 1,\left( {t - 1} \right)^{2}} \right)\),

\(\lbrack T\rbrack_{\mathcal{B}} = \left( {\lbrack 1\rbrack_{\mathcal{B}}\left\lbrack {2t - 2} \right\rbrack_{\mathcal{B}}\left\lbrack {4\left( {t - 1} \right)^{2}} \right\rbrack_{\mathcal{B}}} \right) = \begin{pmatrix}
1 & 0 & 0 \\
0 & 2 & 0 \\
0 & 0 & 4
\end{pmatrix}.\)

This has full rank, so \(T\) is invertible, is an isomorphism, and \(\text{rank}(T) = 3\).

\section{4.3 Exercise 60}

In \(2x_{1} + x_{2} - 2x_{3} = 0\) let

\(\mathcal{A} = \left( {\begin{pmatrix}
1 \\
2 \\
2
\end{pmatrix},\begin{pmatrix}
2 \\
{- 2} \\
1
\end{pmatrix}} \right),\quad\mathcal{B} = \left( {\begin{pmatrix}
1 \\
2 \\
2
\end{pmatrix},\begin{pmatrix}
3 \\
0 \\
3
\end{pmatrix}} \right).\)

The submitted change-of-basis matrices are

\(S_{\mathcal{B}\rightarrow\mathcal{A}} = \begin{pmatrix}
1 & 1 \\
0 & 1
\end{pmatrix},\quad S_{\mathcal{A}\rightarrow\mathcal{B}} = \begin{pmatrix}
1 & {- 1} \\
0 & 1
\end{pmatrix},\)

and \(\left\lbrack {\overrightarrow{a_{1}}\overrightarrow{a_{2}}} \right\rbrack = S_{\mathcal{B}\rightarrow\mathcal{A}}\left\lbrack {\overrightarrow{b_{1}}\overrightarrow{b_{2}}} \right\rbrack\).

\section{5.1 Exercise 6}

For \(u = \begin{pmatrix}
1 \\
{- 1} \\
2 \\
{- 2}
\end{pmatrix}\) and \(v = \begin{pmatrix}
2 \\
3 \\
4 \\
5
\end{pmatrix}\),

\(\cos\theta = \frac{u \cdot v}{\left\| u \right\|\left\| v \right\|} = \frac{- 3}{\sqrt{540}} = \frac{- 1}{2\sqrt{15}}.\)

Thus \(\theta = \arccos\left( {- \frac{\sqrt{15}}{30}} \right) \approx 1.7\) rad (approximately 97.4 degrees).

\section{5.1 Exercise 17}

For \(W = \text{span}\left( {\begin{pmatrix}
1 \\
2 \\
3 \\
4
\end{pmatrix},\begin{pmatrix}
5 \\
6 \\
7 \\
8
\end{pmatrix}} \right)\), the equations for \(x = \begin{pmatrix}
x_{1} \\
x_{2} \\
x_{3} \\
x_{4}
\end{pmatrix} \in W^{\perp}\) reduce as

\(\begin{pmatrix}
1 & 2 & 3 & 4 \\
5 & 6 & 7 & 8
\end{pmatrix}\mapsto\begin{pmatrix}
1 & 0 & {- 1} & {- 2} \\
0 & 1 & 2 & 3
\end{pmatrix}.\)

Thus \(x_{1} = x_{3} + 2x_{4}\), \(x_{2} = - 2x_{3} - 3x_{4}\), and

\(W^{\perp} = \left\{ r\begin{pmatrix}
1 \\
{- 2} \\
1 \\
0
\end{pmatrix} + s\begin{pmatrix}
2 \\
{- 3} \\
0 \\
1
\end{pmatrix}\  \middle| \ r,s \in \mathbb{R} \right\}.\)

\section{5.1 Exercise 26}

With \(u_{1} = \frac{1}{7}\begin{pmatrix}
2 \\
3 \\
6
\end{pmatrix}\) and \(u_{2} = \frac{1}{7}\begin{pmatrix}
3 \\
{- 6} \\
2
\end{pmatrix}\),

\(\text{proj}_{W{(x)}} = \left( {u_{1} \cdot x} \right)u_{1} + \left( {u_{2} \cdot x} \right)u_{2}\)

\(= \left( {2 + 3 + 6} \right)\begin{pmatrix}
2 \\
3 \\
6
\end{pmatrix} + \left( {3 - 6 + 2} \right)\begin{pmatrix}
3 \\
{- 6} \\
2
\end{pmatrix} = \begin{pmatrix}
19 \\
35 \\
64
\end{pmatrix}.\)

\chapter{Part B}

\section{Problem 1}

For ordered bases \(A,B,C\), writing \(C = \left( {c_{1},c_{2},\ldots,c_{n}} \right)\) and taking \(f \in W\),

\(\lbrack f\rbrack_{A} = S_{C\rightarrow A}\lbrack f\rbrack_{C} = S_{B\rightarrow A}\lbrack f\rbrack_{B},\quad\lbrack f\rbrack_{B} = S_{C\rightarrow B}\lbrack f\rbrack_{C}.\)

So

\(S_{C\rightarrow A}\lbrack f\rbrack_{C} = \left( {S_{B\rightarrow A}S_{C\rightarrow B}} \right)\lbrack f\rbrack_{C}.\)

Taking \(\lbrack f\rbrack_{C} = e_{i}\) makes corresponding columns equal; therefore

\(S_{C\rightarrow A} = S_{B\rightarrow A}S_{C\rightarrow B}.\)

Consequently,

\(S_{C\rightarrow A}S_{B\rightarrow C}S_{A\rightarrow B} = S_{B\rightarrow A}S_{C\rightarrow B}S_{B\rightarrow C}S_{A\rightarrow B} = S_{B\rightarrow A}I_{n}S_{A\rightarrow B} = I_{n}.\)

\section{Problem 2}

For \(f_{1} = \sin\left( {2x} \right)\), \(f_{2} = \cos\left( {2x} \right)\), \(f_{3} = e^{3x}\) and \(\mathcal{B} = \left( {f_{1},f_{2},f_{3}} \right)\),

\(\left\lbrack {D\left( f_{1} \right)} \right\rbrack_{\mathcal{B}} = \begin{pmatrix}
0 \\
2 \\
0
\end{pmatrix},\quad\left\lbrack {D\left( f_{2} \right)} \right\rbrack_{\mathcal{B}} = \begin{pmatrix}
{- 2} \\
0 \\
0
\end{pmatrix},\quad\left\lbrack {D\left( f_{3} \right)} \right\rbrack_{\mathcal{B}} = \begin{pmatrix}
0 \\
0 \\
3
\end{pmatrix},\)

so \(\lbrack D\rbrack_{\mathcal{B}} = \begin{pmatrix}
0 & {- 2} & 0 \\
2 & 0 & 0 \\
0 & 0 & 3
\end{pmatrix}\).

The submitted interpretation is that it rotates every vector in \(\mathbb{R}^{2}\) counterclockwise by \(\frac{\pi}{2}\), stretches the \(x,y\) coordinates to twice their length, and the \(z\) coordinate to three times its length.

\section{Problem 3}

Let \(B = \lbrack T\rbrack_{\mathcal{B}}\) and \(C = \lbrack T\rbrack_{\mathcal{C}}\). By the change-of-basis theorem, \(C = S^{- 1}BS\), where \(S = S_{C\rightarrow B}\). The submission proves by induction that

\(B^{k} = S^{- 1}C^{k}S\)

for each integer \(k \geq 1\). The inductive step is

\(B^{k + 1} = B^{k}B = \left( {S^{- 1}C^{k}S} \right)\left( {S^{- 1}CS} \right) = S^{- 1}C^{k + 1}S.\)

Thus \(B^{k}\) and \(C^{k}\) are similar.

For the false kernel claim, it takes

\(A = \begin{pmatrix}
1 & 0 & 0 \\
0 & 1 & 0 \\
0 & 0 & 0
\end{pmatrix},\)

\(\mathcal{E} = \left( {\begin{pmatrix}
1 \\
0 \\
0
\end{pmatrix},\begin{pmatrix}
0 \\
1 \\
0
\end{pmatrix},\begin{pmatrix}
0 \\
0 \\
1
\end{pmatrix}} \right),\)

\(\mathcal{B} = \left( {\begin{pmatrix}
0 \\
1 \\
0
\end{pmatrix},\begin{pmatrix}
1 \\
0 \\
0
\end{pmatrix},\begin{pmatrix}
0 \\
0 \\
1
\end{pmatrix}} \right).\)

It records

\(\lbrack A\rbrack_{\mathcal{E}} = A,\quad\lbrack A\rbrack_{\mathcal{B}} = \begin{pmatrix}
1 & 0 & 0 \\
0 & 0 & 0 \\
0 & 1 & 1
\end{pmatrix},\)

but

\(\text{ker}\left( \lbrack A\rbrack_{\mathcal{E}} \right) = \left\{ r\begin{pmatrix}
0 \\
0 \\
1
\end{pmatrix}\  \middle| \ r \in \mathbb{R} \right\},\)

\(\text{ker}\left( \lbrack A\rbrack_{\mathcal{B}} \right) = \left\{ r\begin{pmatrix}
0 \\
1 \\
{- 1}
\end{pmatrix}\  \middle| \ r \in \mathbb{R} \right\}.\)

For equal nullities, rank-nullity gives

\(\text{dim}\left( {\text{ker}\ B} \right) = n - \text{rank}(B),\quad\text{dim}\left( {\text{ker}\ C} \right) = n - \text{rank}(C).\)

The page proves \(\text{rank}\left( {MN} \right) \leq \text{rank}(N)\) from \(\text{ker}\ N \subseteq \text{ker}\left( {MN} \right)\), the analogous bound for \(M\) by transposition, and equality \(\text{rank}\left( {MN} \right) = \text{rank}(N)\) when \(M\) is invertible. Since \(B = S^{- 1}CS\), \(\text{rank}(B) = \text{rank}(C)\), so the kernel dimensions agree.

\section{Problem 4}

For \(T:U\rightarrow W\), bases \(\mathcal{B} = \left( {u_{1},\ldots,u_{k}} \right)\) and \(\mathcal{C} = \left( {w_{1},\ldots,w_{d}} \right)\), define

\(T'\left( \lbrack u\rbrack_{\mathcal{B}} \right) = \lbrack w\rbrack_{\mathcal{C}}\quad\text{whenever}\quad T(u) = w.\)

The source diagram verifies

\(T' \circ L_{\mathcal{B}}(u) = T'\left( \lbrack u\rbrack_{\mathcal{B}} \right) = \lbrack w\rbrack_{\mathcal{C}} = L_{\mathcal{C}} \circ T(u),\)

hence \(T' \circ L_{\mathcal{B}} = L_{\mathcal{C}} \circ T\). Thus

\(\left\lbrack {T(u)} \right\rbrack_{\mathcal{C}} = \lbrack T\rbrack_{\mathcal{B},\mathcal{C}}\lbrack u\rbrack_{\mathcal{B}}.\)

For \(u = a_{1}u_{1} + \ldots + a_{k}u_{k}\),

\(\left\lbrack {T(u)} \right\rbrack_{\mathcal{C}} = a_{1}\left\lbrack {T\left( u_{1} \right)} \right\rbrack_{\mathcal{C}} + \ldots + a_{{k{\lbrack{T{(u_{k})}}\rbrack}}_{\mathcal{C}}},\)

and therefore

\(\lbrack T\rbrack_{\mathcal{B},\mathcal{C}} = \left( {\left\lbrack {T\left( u_{1} \right)} \right\rbrack_{\mathcal{C}}\left\lbrack {T\left( u_{2} \right)} \right\rbrack_{\mathcal{C}}\ldots\left\lbrack {T\left( u_{k} \right)} \right\rbrack_{\mathcal{C}}} \right).\)

\section{Problem 5}

For \(f_{1} = \sin x\), \(f_{2} = \cos x\), \(f_{3} = e^{x}\):

\(T\left( f_{1} \right) = x - \frac{x^{3}}{6},\quad T\left( f_{2} \right) = 1 - \frac{x^{2}}{2},\)

\(T\left( f_{3} \right) = 1 + x + \frac{x^{2}}{2} + \frac{x^{3}}{6}.\)

The source chooses \(\mathcal{C} = \left( {1,x,\frac{x^{2}}{2},\frac{x^{3}}{6}} \right)\). For \(\mathcal{B} = \left( {f_{1} + f_{2},f_{1} - f_{2},f_{3} + f_{1}} \right)\) it computes

\(\left\lbrack {T\left( {f_{1} + f_{2}} \right)} \right\rbrack_{\mathcal{C}} = \begin{pmatrix}
1 \\
1 \\
{- 1} \\
{- 1}
\end{pmatrix},\)

\(\left\lbrack {T\left( {f_{1} - f_{2}} \right)} \right\rbrack_{\mathcal{C}} = \begin{pmatrix}
{- 1} \\
1 \\
1 \\
{- 1}
\end{pmatrix},\)

\(\left\lbrack {T\left( {f_{3} + f_{1}} \right)} \right\rbrack_{\mathcal{C}} = \begin{pmatrix}
1 \\
2 \\
1 \\
0
\end{pmatrix},\)

so

\(\lbrack T\rbrack_{\mathcal{B},\mathcal{C}} = \begin{pmatrix}
1 & {- 1} & 1 \\
1 & 1 & 2 \\
{- 1} & 1 & 1 \\
{- 1} & {- 1} & 0
\end{pmatrix}.\)

\section{Problem 6}

Let \(A = \begin{pmatrix}
{- 6} & {- 30} \\
{- 30} & 19
\end{pmatrix}\) and \(V = \text{span}\left( \begin{pmatrix}
3 \\
2
\end{pmatrix} \right)\). For \(v = a\begin{pmatrix}
3 \\
2
\end{pmatrix}\),

\(Av = a\begin{pmatrix}
{- 78} \\
{- 52}
\end{pmatrix} = - 26a\begin{pmatrix}
3 \\
2
\end{pmatrix} \in V.\)

\(V^{\perp} = \left\{ r\begin{pmatrix}
{- 2} \\
3
\end{pmatrix}\  \middle| \ r \in \mathbb{R} \right\}\), and

\(A\begin{pmatrix}
{- 2r} \\
{3r}
\end{pmatrix} = r\begin{pmatrix}
{- 78} \\
117
\end{pmatrix} = - 39r\begin{pmatrix}
{- 2} \\
3
\end{pmatrix} \in V^{\perp}.\)

With \(\mathcal{B} = \left( {\begin{pmatrix}
{- 2} \\
3
\end{pmatrix},\begin{pmatrix}
3 \\
2
\end{pmatrix}} \right)\),

\(\lbrack T\rbrack_{\mathcal{B}} = \begin{pmatrix}
26 & 0 \\
0 & {- 39}
\end{pmatrix},\quad\left\lbrack T^{10} \right\rbrack_{\mathcal{B}} = \begin{pmatrix}
26^{10} & 0 \\
0 & 39^{10}
\end{pmatrix}.\)

The source uses

\(S_{\mathcal{B}\rightarrow\mathcal{E}} = \begin{pmatrix}
{- 2} & 3 \\
3 & 2
\end{pmatrix},\quad S_{\mathcal{B}\rightarrow\mathcal{E}}^{- 1} = \begin{pmatrix}
\frac{2}{13} & {- \frac{3}{13}} \\
{- \frac{3}{13}} & {- \frac{2}{13}}
\end{pmatrix}\)

to obtain

\(\left\lbrack T^{10} \right\rbrack_{\mathcal{E}} = \begin{pmatrix}
\frac{4 \cdot 26^{10} + 9 \cdot 39^{10}}{13} & \frac{- 6 \cdot 26^{10} + 6 \cdot 39^{10}}{13} \\
\frac{- 6 \cdot 26^{10} + 6 \cdot 39^{10}}{13} & \frac{9 \cdot 26^{10} + 4 \cdot 39^{10}}{13}
\end{pmatrix}.\)
